\chapter{Tableaux complémentaires}
\label{chap:B-tableaux}

Cette annexe donne les tableaux qui ne figurent pas dans les chapitres, puis un index de tous les
tableaux générés et la liste des fichiers de résultats d'où ils viennent. Chaque tableau généré est
écrit par \file{report/tools/tables.py} à partir d'un fichier JSON de résultats : aucun nombre n'y
est tapé à la main.

\section{Le RL de grands modèles de langue dans le sondage}
\label{sec:b-lrl}

Le premier tirage de 2026 compte 27 articles qui appliquent ce qu'ils appellent RL à de grands
modèles de langue (famille X5 \texttt{llm}, \cref{chap:A-grille}). Les verdicts du sondage valent pour
le RL classique. Le \cref{tab:sondage-lrl} met ces articles en regard du RL classique de 2026, critère
par critère, là où le fichier de résultats le permet. Les 18 articles \texttt{llm} lus en chemin dans
la file de l'élargissement n'y entrent pas : le tri sur le titre les avait filtrés, si bien qu'ils ne
forment pas un échantillon aléatoire (écart 12 du sondage). Le tableau se lit comme une description,
sans verdict : il compare deux familles d'articles, pas deux époques.

\section{Index des tableaux générés}
\label{sec:b-index}

\begin{small}
\begin{longtable}{@{}p{0.19\linewidth}p{0.43\linewidth}p{0.3\linewidth}@{}}
\caption{Où trouver chaque tableau généré.}\label{tab:b-index}\\
\toprule
Tableau & Contenu & Fichier source \\
\midrule
\endfirsthead
\toprule
Tableau & Contenu & Fichier source \\
\midrule
\endhead
\bottomrule
\endfoot
\Cref{tab:sondage-pineau} & Critères de Pineau, 2018 contre RL classique de 2026 &
  \file{extension\_stats.json}, \file{results\_final.json} \\
\Cref{tab:sondage-extension} & Nos critères Y1 à Y12, mêmes colonnes & \file{extension\_stats.json} \\
\Cref{tab:sondage-lrl} & RL de grands modèles contre RL classique, 2026 & \file{extension\_stats.json} \\
\Cref{tab:trajectoire} & Part par pratique et par année, test de tendance & \file{trajectory.json} \\
\Cref{tab:accord} & Accord des doubles lectures à l'aveugle & \file{agree\_v2.json} \\
\Cref{tab:robustesse} & Contrôles de robustesse du sondage & \file{robustness.json} et variantes \\
\Cref{tab:e1} & E1 : trente essais de PPO, faux positifs & \file{e1.json} \\
\Cref{tab:b1} & B1 : le $t$ de 9,09, deux façons de le calculer & \file{b1\_script.json}, \file{b1\_trpo.json} \\
\Cref{tab:e2} & E2 : classement par tâche, dix et vingt graines aléatoires, budget égal & \file{e2*.json} \\
\Cref{tab:e2-bonferroni} & E2 avec intervalles de Bonferroni & \file{e2\_bonferroni*.json} \\
\Cref{tab:e3} & E3 : PPO PyTorch contre PPO JAX & \file{e3*.json} \\
\Cref{tab:e4} & E4 : les trois meilleurs sur dix & \file{e4.json}, \file{e4\_r10.json} \\
\Cref{tab:e5} & E5 : combien de graines aléatoires retrouvent l'ordre & \file{e5*.json} \\
\Cref{tab:e6} & E6 : la référence réglée & \file{e6*.json} \\
\Cref{tab:e7} & E7 : score gardé sous physique modifiée & \file{e7*.json} \\
\Cref{tab:dose} & Physique modifiée par degrés & \file{dose*.json} \\
\Cref{tab:b3} & B3 : la mesure change-t-elle le classement & \file{b3*.json} \\
\Cref{tab:grpo} & GRPO : chaque essai & \file{summary\_5.json}, \file{gains\_5.json}, \file{extra\_5.json} \\
\Cref{tab:ecarts} & Les écarts des trois protocoles & les trois protocoles \\
\Cref{tab:env} & Versions des environnements & \file{envs/*.lock} \\
\Cref{tab:calcul} & Le calcul par programme et par tâche & \file{sim/logs}, \file{sim/results} \\
\Cref{tab:articles-2026} & Les articles de 2026 lus & \file{final\_2026.json} et file A3 \\
\Cref{tab:articles-2018} & Les articles de 2018 lus & \file{calib\_2018\_v2.json} \\
\end{longtable}
\end{small}

\section{Les fichiers de résultats}
\label{sec:b-fichiers}

Chaque fichier JSON de résultats vient d'un seul script, relancé par \texttt{make analysis}
(\cref{sec:make}). Les détails critère par critère et paire par paire que les tableaux résument se
trouvent dans ces fichiers.

\begin{small}
\begin{longtable}{@{}p{0.3\linewidth}p{0.64\linewidth}@{}}
\caption{Les scripts d'analyse et ce que contient leur fichier JSON.}\label{tab:b-fichiers}\\
\toprule
Script & Contenu du fichier de résultats \\
\midrule
\endfirsthead
\toprule
Script & Contenu du fichier de résultats \\
\midrule
\endhead
\bottomrule
\endfoot
\multicolumn{2}{@{}l}{\textit{Sondage, dossier \file{study/}}}\\
\file{analyze.py} & Parts du premier tirage de 2026 (\file{results\_2026.json}). \\
\file{analyze\_final.py} & Parts de 2026 et de la calibration de 2018 depuis les lectures brutes et
  l'harmonisation (\file{results\_final.json}). \\
\file{audit\_stats.py} & Audit statistique de la comparaison entre 2018 et 2026 : différences de
  Newcombe, Fisher, Holm, accord des lecteurs (\file{audit\_stats.json}). \\
\file{agree\_v2.py} & Accord entre lecteurs sur les lectures brutes, avant harmonisation
  (\file{agree\_v2.json}). \\
\file{extension\_stats.py} & Comparaison principale : RL classique de 2026 contre notre lecture de
  2018, critère par critère, avec Wilson, Newcombe, Fisher et Holm (\file{extension\_stats.json} ;
  variantes sans le modèle voisin et avec les décisions retournées). \\
\file{robustness.py} & Coupes de robustesse et test de la thèse du coût (\file{robustness.json}). \\
\file{title\_filter.py}, \file{title\_filter\_compare.py} & Le filtre de titre : part de RL caché
  et comparaison des vingt articles cachés (\file{title\_filter*.json}). \\
\file{trajectory.py} & La trajectoire de 2018 à 2026 : parts par année, Cochran-Armitage, écarts
  immédiat et tardif (\file{trajectory.json}). \\
\file{falsify\_no.py} & La recherche contre les décisions \enquote{non} (\file{falsify\_no.json}). \\
\file{checklist.py} & Les listes de contrôle déclarées à NeurIPS 2025 contre nos lectures
  (\file{checklist.json}). \\
\multicolumn{2}{@{}l}{\textit{Simulations, dossier \file{sim/analysis/}}}\\
\file{e1.py}, \file{e1\_curve.py} & E1 : trente essais de PPO, test de Welch sur les scores finaux
  puis sur les courbes. \\
\file{b1\_trpo.py}, \file{b1\_script.py} & B1 : les lectures du $t$ moyen de 2018 sur TRPO et PPO, puis
  le calcul du script de Henderson. \\
\file{e2.py}, \file{e2\_bonferroni.py} & E2 : moyenne interquartile par tâche, probabilités
  d'amélioration, changements de classement, puis les mêmes intervalles au niveau simultané de
  Bonferroni. \\
\file{e3.py} & E3 : les deux implémentations de PPO. \\
\file{e4.py} & E4 : moyenne et dispersion des dix essais puis des trois meilleurs, renversements. \\
\file{e5.py} & E5 : part des tirages qui retrouvent l'ordre, pour 3, 5, 10 et 20 graines aléatoires. \\
\file{e6.py}, \file{e6\_diag.py}, \file{e6\_official.py} & E6 : défaut contre Zoo, diagnostic de
  la recette, politiques officielles évaluées avec notre protocole. \\
\file{e7.py}, \file{e7\_rank.py} & E7 : part du score gardée sous physique modifiée, puis le classement
  descriptif. \\
\file{dose.py} & La physique modifiée par degrés. \\
\file{b3.py} & B3 : trois mesures par essai et les renversements entre mesures. \\
\multicolumn{2}{@{}l}{\textit{Volet GRPO, dossier \file{sim/llm/}}}\\
\file{summary.py} & Pass@1 aux pas 0 et 100 de chaque essai, gains, intervalles, sensibilité au
  gain manquant, essais CPU seuls (\file{summary*.json}, \file{gains*.json}, \file{extra*.json}). \\
\end{longtable}
\end{small}

Les suffixes des fichiers des simulations suivent l'écart 12 : \file{\_s20} pour vingt graines aléatoires,
\file{\_r10} pour la réplication de E4 sur les graines aléatoires 10 à 19, \file{\_equal} et
\file{\_equal\_s20} pour le budget égal. Les fichiers du volet GRPO suivent l'écart 14 : sans
suffixe pour trois graines aléatoires, \file{\_5} pour cinq.
